\chapter{机器学习流程}
\section{数据处理}
\subsection{问题定义与规划}
要解决什么问题？ 本研究旨在解决复杂水下物理环境（如极化带效应、涡流效应）中，对目标产生的非平稳、稀疏电荷信号进行实时探测、特征提取与多目标识别的问题。核心挑战在于从强背景噪声中有效分离弱目标信号，并实现对目标电荷时空分布特性的准确建模与识别。

如何将其转化为机器学习问题？ 一个多任务学习问题，包含：

信号检测：二分类问题（有目标/无目标）；特征提取：时间序列特征工程与表示学习；多目标识别：多分类或序列标注问题（若目标类型多样）。

成功的标准是什么？ 检测准确率 > 95\%，误报率 < 5\%；目标识别准确率 > 90\%；系统响应延迟低于 1 秒，满足实时性要求；在仿真与实测数据中均保持一致性表现。

需要哪些数据？ 16 通道电压时间序列数据（采样频率需明确）；标注数据：包括目标出现时间、类型、位置（若可行）；背景噪声数据（无目标时段）；仿真数据（COMSOL 有限元模拟结果）；实际实验数据（水池或海域实测）。

\subsection{数据收集与准备}
数据收集：通过 MEMS 传感器阵列采集 16 通道电压时序数据，每段持续约 20 秒；同步记录环境参数（水温、流速等）与目标状态（类型、运动轨迹）。

数据清洗：处理缺失值：采用插值或前后帧填充；异常值处理：基于统计方法（如 3σ 准则）识别并修正或剔除；去除重复采集数据。

数据探索与可视化：绘制各通道时域波形、频谱图、时频图（如 STFT 或小波变换）；分析通道间相关性、背景噪声分布、目标信号典型模式。

数据标注：对实测数据人工标注目标出现区间与类型；仿真数据自带标签（目标类型、电荷分布、运动状态）。

\subsection{特征工程}
特征构造：时域特征：均值、方差、峰值、过零率、短时能量等；频域特征：FFT 幅度谱、相位谱、频谱重心、带宽；时频特征：小波系数、STFT 频谱图；空间特征：通道间差分、协方差矩阵、阵列响应模式。

特征转换：对数据进行标准化、归一化、对数变换等，使其更适合模型处理。

特征选择：使用互信息、卡方检验、LASSO 等方法筛选关键特征；考虑使用遗传算法进行特征子集优化。

\subsection{模型选择与训练}
模型选择：基线模型：SVM、随机森林、LSTM、1D-CNN；

进阶模型：注意力机制+CNN/LSTM、Transformer 时序模型；多任务学习框架：联合 训练检测与识别任务。

算法库包括：线性模型、决策树、神经网络、支持向量机（SVM）、聚类算法（如K-Means）等。

模型训练：

训练方法：监督学习：使用标注数据训练；训练/验证集划分：按时间序列分段，避免随机划分导致信息泄漏

优化算法：梯度下降是训练神经网络最常用的方法。Adam、SGD；可尝试遗传算法进行超参数搜索。

\subsection{模型评估与调优}
模型评估：评估指标：准确率、精确率、召回率、F1-score、AUC（检测任务）；分类准确率、混淆矩阵（识别任务）；测试集：完全独立于训练集的实测数据。

模型调优：超参数调优：网格搜索、贝叶斯优化或遗传算法；过拟合控制：Dropout、权重正则化、早停法；数据增强：添加噪声、时域拉伸、通道抖动。

\subsection{模型部署与监控}
模型通过评估后，工作并未结束。必须将其投入使用并持续维护。

模型部署：将训练好的模型集成到现有的生产环境中，提供API接口、集成到应用程序中等，使其能够对真实数据做出预测。

模型监控：持续监控模型在生产环境中的性能。因为现实世界的数据分布可能会随时间发生变化（概念漂移），导致模型性能下降。

模型更新与维护：定期用新的数据重新训练模型（迭代过程），以保持其预测的准确性。

% TODO: 草稿中后续重复出现了一段通用机器学习流程模板，已不重复迁移。必要时可与本章内容合并修订。
